%-------------------------
% Resume in Latex
% Author : Jake Gutierrez
% Based off of: https://github.com/sb2nov/resume
% License : MIT
%------------------------

\documentclass[a4paper,11pt]{article}

\usepackage{latexsym}

\usepackage{titlesec}
\usepackage{marvosym}
\usepackage[usenames,dvipsnames]{color}
\usepackage{verbatim}
\usepackage{enumitem}
\usepackage[hidelinks]{hyperref}
\usepackage{fancyhdr}
\usepackage[UTF8,scheme=plain,fontset=fandol]{ctex}
\usepackage{tabularx}
\usepackage{iftex}
\ifPDFTeX
  \input{glyphtounicode}
\else
  \usepackage{fontspec}
  \setmainfont{texgyretermes-regular.otf}[BoldFont=texgyretermes-bold.otf,ItalicFont=texgyretermes-italic.otf,BoldItalicFont=texgyretermes-bolditalic.otf]
\fi


%----------FONT OPTIONS----------
% sans-serif
% \usepackage[sfdefault]{FiraSans}
% \usepackage[sfdefault]{roboto}
% \usepackage[sfdefault]{noto-sans}
% \usepackage[default]{sourcesanspro}

% serif
% \usepackage{CormorantGaramond}
% \usepackage{charter}


\pagestyle{fancy}
\fancyhf{} % clear all header and footer fields
\fancyfoot{}
\renewcommand{\headrulewidth}{0pt}
\renewcommand{\footrulewidth}{0pt}

\usepackage[a4paper,top=15mm,bottom=15mm,left=18mm,right=18mm]{geometry}

\urlstyle{same}

\raggedbottom
\raggedright
\setlength{\tabcolsep}{0in}

% Sections formatting
\titleformat{\section}{
  \vspace{-4pt}\scshape\raggedright\large
}{}{0em}{}[\color{black}\titlerule \vspace{-5pt}]

% Ensure that generate pdf is machine readable/ATS parsable
\ifPDFTeX\pdfgentounicode=1\fi

%-------------------------
% Custom commands
\newcommand{\resumeItem}[1]{
  \item\small{
    {#1 \vspace{-2pt}}
  }
}

\newcommand{\resumeSubheading}[4]{
  \vspace{-2pt}\item
    \begin{tabular*}{0.97\textwidth}[t]{l@{\extracolsep{\fill}}r}
      \textbf{#1} & #2 \\
      \textit{\small#3} & \textit{\small #4} \\
    \end{tabular*}\vspace{-1pt}
}

\newcommand{\resumeSubSubheading}[2]{
    \item
    \begin{tabular*}{0.97\textwidth}{l@{\extracolsep{\fill}}r}
      \textit{\small#1} & \textit{\small #2} \\
    \end{tabular*}\vspace{-1pt}
}

\newcommand{\resumeProjectHeading}[2]{
    \item
    \begin{tabular*}{0.97\textwidth}{l@{\extracolsep{\fill}}r}
      \small#1 & #2 \\
    \end{tabular*}\vspace{-1pt}
}

\newcommand{\resumeSubItem}[1]{\resumeItem{#1}\vspace{-4pt}}

\renewcommand\labelitemii{$\vcenter{\hbox{\tiny$\bullet$}}$}

\newcommand{\resumeSubHeadingListStart}{\begin{itemize}[leftmargin=0.15in,label={},topsep=2pt,itemsep=2pt,parsep=0pt]}
\newcommand{\resumeSubHeadingListEnd}{\end{itemize}}
\newcommand{\resumeItemListStart}{\begin{itemize}}
\newcommand{\resumeItemListEnd}{\end{itemize}\vspace{-5pt}}

%-------------------------------------------
%%%%%%  RESUME STARTS HERE  %%%%%%%%%%%%%%%%%%%%%%%%%%%%



% Chinese adaptation: portable Fandol fonts included in TeX Live.
\renewcommand{\baselinestretch}{1.04}
\titleformat{\section}{\vspace{-2pt}\bfseries\large}{}{0em}{}[\color{black}\titlerule\vspace{-4pt}]

% Compile with XeLaTeX. Edit the content below to update your information.

% Research entries: full titles, publication status, and two labeled paragraphs.
\newcommand{\researchEntry}[4]{%
  \par\addvspace{10pt}\noindent
  {\normalsize\bfseries #1\par}\nopagebreak
  {\footnotesize\itshape #2\par}\nopagebreak\vspace{2pt}
  {\small\noindent\textbf{研究内容：}#3\par}\nopagebreak\vspace{2pt}
  {\small\noindent\textbf{个人贡献：}#4\par}
}
\begin{document}
\begin{center}
\textbf{\Huge 晏景豪} \\[3pt]
\small 中国科学技术大学 · 博士研究生 \\[2pt]
\small 具身智能   |   灵巧手精细操作   |   多模态感知融合 \\[2pt]
\small 邮箱：\href{mailto:jinghaoyan943@gmail.com}{jinghaoyan943@gmail.com} \\[2pt]
\end{center}

\section{教育背景}
\resumeSubHeadingListStart
\resumeSubheading{中国科学技术大学}{2026.09 -- 至今}{自动化系 · 电子信息专业 · 博士研究生}{}
\resumeSubHeadingListEnd
\resumeSubHeadingListStart
\resumeSubheading{同济大学}{2022.09 -- 2026.06}{计算机科学与技术专业 · 本科}{}
\item\small 平均绩点：4.97/5.00   |   专业排名：1/118（前 0.85\%）
\resumeSubHeadingListEnd

\section{科研与实习经历}
\noindent\textbf{上海人工智能实验室}\hfill 2025.10 -- 至今\par
{\small 科研实习生\par}
\researchEntry{MegaParts: Scaling Part-Aware 3D Object Generation to 300 Parts via Token-Efficient Autoregressive Modeling}{SIGGRAPH Asia 2026 已接收}{面向复杂多部件物体生成，结合自适应长度的离散几何编码与统一自回归序列建模，联合生成物体及部件包围盒与部件几何。通过压缩几何 token 降低长序列开销，将生成规模扩展至 300 个部件，同时保留精细结构与部件级控制能力。}{负责基线方法的实验运行、结果统计与对比分析，为生成质量与方法有效性评估提供实验支持；熟悉项目整体框架、部件表示方式及实验设计。}
\researchEntry{Code-as-Room: Generating 3D Rooms from Top-Down View Images via Agentic Code Synthesis}{预印本}{面向俯视图到三维室内场景的生成，利用多模态大模型解析场景元素与空间关系，以 Blender 代码表示几何、材质和光照。通过分阶段执行与跨阶段记忆，缓解 Agent 在完整场景生成中的上下文遗忘、执行不稳定与循环问题。}{负责三维生成模型在 Agent 系统中的集成，使 Agent 能够调用生成模型制作难以直接通过代码构建的物体，补充代码建模能力并扩展可生成物体的范围。}
\researchEntry{From Parts to Sim: Bottom-Up Autoregressive Generation of Sim-Ready 3D Assets}{在审}{面向可交互、可仿真的三维资产生成，以自底向上的自回归框架关联精细几何部件与刚体运动学结构，联合建模几何、关节及物理属性。生成资产支持 USD/URDF 导出，并用于仿真交互轨迹采集与机器人策略的仿真到现实迁移。}{负责基线实验运行、结果统计与分析，以及材质生成管线搭建：结合部件 UV 图集、物体级材质生成与部件材质重映射，在保留关节结构的同时维持跨部件外观一致性。}

\section{荣誉奖励}
\resumeSubHeadingListStart
\item\small 国家奖学金（Top 3\%）\hfill 2022--2023
\item\small 国家奖学金（Top 3\%）\hfill 2023--2024
\item\small 国家奖学金（Top 3\%）\hfill 2024--2025
\item\small 上海市优秀毕业生（Top 2\%）\hfill 2026
\resumeSubHeadingListEnd

\section{专业技能}
\begin{itemize}[leftmargin=0.15in,label={},itemsep=1pt]
\item\small\textbf{编程与框架：} Python、PyTorch
\item\small\textbf{深度学习：} 熟悉模型部署与实验评估
\item\small\textbf{三维工具：} 具备 Blender 与三维资产处理经验
\end{itemize}
\end{document}
